\section{Introduction}
Post-training has become a central stage in the development of modern large language models, transforming a pre-trained model into one that can reliably follow instructions~\citep{ouyang2022training}, reason over complex inputs~\citep{deep-seek-r1}, and interact with users across a broad range of applications~\citep{openclaw-rl2026,yue2026interactivebenchmarks}. In practice, this stage is rarely optimized for a single task or capability. Instead, post-training data typically mixes many heterogeneous tasks, including mathematical reasoning, coding, factual question answering, instruction following, tool use, and other domain-specific problems, often within the mini-batch or the same training run. As a result, post-training is naturally a multi-task learning problem~\citep{nvidia2025nemotron3nanoopen}. Under a finite training budget, choosing this mixture becomes especially important, as tasks can differ substantially in difficulty and learning dynamics, causing some tasks to improve faster than others under the same allocation~\citep{wu2025imbalancedgradientsrlposttraining}. Furthermore, gradients from different tasks can interfere through shared model updates, suppressing useful task-specific learning signals and reducing overall training efficiency~\citep{pcgrad2020,cagrad2021}.

Existing multi-task learning methods address these issues through two broad pathways: gradient-surgery methods that explicitly compute and manipulate task-wise gradients~\citep{pcgrad2020,cagrad2021}, or fairness and multi-objective-aware methods that reweigh task losses to favor those making the least amount of progress~\citep{cagrad2021,sener2018mgda}. While these approaches provide principled mechanisms for addressing task interference or balancing objectives, neither pathway can be relied upon directly for LLM post-training. Gradient-based methods such as PCGrad~\citep{pcgrad2020} require computing and storing task-wise gradients before constructing a joint update, introducing substantial additional backward-pass and memory costs for models with billions of parameters. These costs scale further with the number of tasks linearly and are difficult to reconcile with highly optimized LLM training pipelines based on gradient accumulation, parameter sharding, and distributed optimization~\citep{sheng2024hybridflow}. Fairness- and multi-objective-aware methods such as MGDA instead focus on convergence toward balanced or Pareto-stationary solutions. However, LLM post-training is typically conducted under a fixed compute budget and stopped well before convergence, making finite-horizon training efficiency more relevant than asymptotic behavior. These limitations make both classes of methods poorly matched to the practical requirements of contemporary LLM post-training.

In this paper, we formulate multi-task post-training as an online convex optimization problem in which the only control variable is the task quota assigned to each mini-batch. Across training steps, the allocation policy seeks to maximize cumulative policy improvement under a set of ``\textit{no-task-left-behind}'' constraints. These constraints measure, for each task, how much the policy improvement induced by the current allocation can lag behind a counterfactual uniform allocation. To make the constrained problem tractable, we introduce a directional concavity assumption on the task-improvement surface, which linearizes the constraints. The resulting formulation leads to a primal-dual mirror-descent-like algorithm, which we dub \textbf{Primal-Dual Curvature Matching (PDCM)}, that dynamically adjusts task quotas using first-order information about how policy improvement changes with allocation.

The closest prior formulation, Aioli~\citep{chen2025aioli}, casts pre-training data mixing as the same kind of optimization problem but closes it with a power mixing law tailored to pre-training dynamics (\Cref{sec:bg_constrained_allocation}). PDCM assumes only directional concavity of the task-improvement and utilizes local first-order information without over-prescribing the general shape of learning curve.

\textcolor{red}{We need to add a paragraph about the empirical performance of PDCM, as the results come in.}

Our contributions are summarized as follows:
\begin{itemize}
\item We formulate multi-task post-training as a constrained optimization problem over task allocation, maximizing cumulative policy improvement while requiring each task to remain close to its per-step policy improvement under a uniform-allocation counterfactual; we further derive a tractable formulation by linearizing these constraints under a directional concavity assumption based on previous empirical works~\citep{ye2025datamixinglawsoptimizing,chen2025aioli,kang2025autoscalescaleawaredatamixing}.

\item We introduce Primal-Dual Curvature Matching (PDCM), a \textbf{lightweight task-allocation} algorithm that dynamically adjusts per-task training quotas using observed marginal policy improvement and task-wise constraint signals. We provide theoretical guarantees on both cumulative utility and constraint violation under the resulting primal-dual updates.

\item We evaluate PDCM across both \textbf{agentic and non-agentic} multi-task post-training settings and show that dynamic task allocation improves finite-budget training efficiency over static task mixtures, gradient-surgery methods, and multi-objective methods with asymptotic guarantees, while maintaining balanced progress across tasks. These gains come with negligible computational overhead relative to the underlying training procedure. Furthermore, we show that PDCM, which relies only on a non-parametric concavity assumption on the task improvement surface, consistently outperforms the existing data-allocation method~\citep{chen2025aioli} which rely on stronger parametric assumptions.

\end{itemize}